\chapter{The Prize, the Bill, and the Third Path}
\label{c:science}

\begin{glance}
The largest value in AI is not code, documents or enterprise workflow but the acceleration of science, because everything else redistributes an existing surplus and science creates new ones. That turns the price of frontier tokens into a tax on the one activity whose returns compound into everything else. A measurement-anchored sizing framework finds that announced national research infrastructures provide a tenth to a fiftieth of what their scientific populations need, and that a Regime~S workload share of one seventh---a seventh of \emph{scientific} compute, which is a different claim from any share in general-consumer or enterprise AI---becomes a budget share of three quarters to nine tenths once purchased tokens carry a twenty-to-sixty-fold premium. For every country that is neither the United States nor China both available strategies fail---dependence on foreign providers, or an isolated national program at sub-scale---and the third path is an international, institutionally neutral consortium: federated compute, scientific data stewardship, open pretraining as insurance, models with lifecycle obligations, and a certification authority neither bloc controls. It is the constructive branch of an otherwise pessimistic book, and, if it succeeds, a partial falsifier of its strongest claim.
\end{glance}

\section{The prize}

What AI for science has delivered follows one pattern: an AI system replaced a physics-simulation-bounded computation and was validated against an established operational baseline. Weather is the strongest case---the European centre's AI forecast operational since February 2025 at about a thousandth of the energy per forecast~\cite{ecmwf-aifs}, the American service's sixteen-day forecast in forty minutes at 0.3\% of conventional compute. Structural biology is the most diffused---more than 200 million predicted structures, three million users, a Nobel Prize, and users submitting more than 40\% more novel experimental structures~\cite{alphafold-impact}. Mathematics produced the first credible autonomous results, and one drug candidate discovered end to end by AI reached a controlled clinical readout~\cite{rentosertib}. The failure ledger is printed with equal weight: the most-cited economic study of AI in science was fabricated and repudiated~\cite{mit-repudiation}; materials discovery has not survived validation intact; the European centre discontinued all four external AI weather models it had been running~\cite{ecmwf-farewell}; autonomous discovery agents produce mostly flawed candidates; and the journal pipeline is being damaged by AI-generated submissions. It delivers where the task is a well-posed prediction problem with an expensive physics-based incumbent and an established baseline, and has so far mostly failed where the task is judgement.

\emph{October 2026: output is not discovery.} On 6~October OpenAI released 722 mathematics manuscripts from an ``unreleased internal model'' it has not named, and this book attributes them to no named model~\cite{sci-rw-openai} [P]; it withdrew three and revised fourteen within a day [A], a Lean formalization of its Navier--Stokes claim was shown to prove weaker statements than the paper~\cite{sci-navier-lean} [A], and its 26~October claims are unverified~\cite{sci-openai-oct26} [P]. The month's one peer-reviewed result, AlphaProof Nexus, solved 9 of 353 Erd\H{o}s problems, and a plain model with Lean feedback solved all nine~\cite{sci-alphaproof-nexus} [V]: verification is the binding stage. No 2026 Nobel prize had an AI component~\cite{sci-nobel-2026} [V] \full{24}.

The prize is nevertheless the largest because of four numbers. The denominator: global R\&D is about \$3.1 trillion a year at purchasing-power parity, twice the software market, and frontier-laboratory revenue is under 2\% of it~\cite{nsf-indicators}. The return: \$4--13 of social benefit in present value per dollar of R\&D~\cite{jones-summers}. The bottleneck: US research productivity has fallen roughly forty-fold since the 1930s, offset only by twenty-three times more effort~\cite{ideas-harder}, and the published macroeconomic estimates for AI---from a fraction of a percent to thirty percent of growth---are separated entirely by whether AI is assumed to automate innovation itself; the most-cited skeptical estimate excludes protein-folding-type advances as unlikely within ten years, and that exclusion is the entire disagreement, because structure prediction, operational AI weather and designed enzymes all arrived inside the window~\cite{acemoglu}. And the flywheel: AI for science is the only application that feeds back into its own means of production. The book replaces the customary ``3--30\%'' with a value-conversion funnel whose every arrow has an efficiency below one; task-weighted, the central estimate is an uplift of 10.74\% of research output, with a median of 11.61\% and a model-conditional 10th-to-90th-percentile range of 6.73--18.98\% [M]---from a Monte Carlo that ships with the manuscript rather than being promised, 200,000 draws at a fixed seed over triangular marginals with correlated draws, so the interval is reproducible and is conditional on assumed distributions rather than sampled from the world. Even the all-pessimistic corner, 0.83\%, is a flow of \$25.7 billion a year, and the tenth percentile \$209 billion; the present value of either exceeds the cost of national infrastructure by a large multiple. Two consequences carry into the strategy. A scientific result produced by an unverifiable artifact through an unauditable chain is not a scientific result, which makes certification a precondition of the prize rather than an accessory to it. And the workload of scientific AI is long-context agentic orchestration over large batches of scientific execution---precisely the architecture the compute chapter described. China leads the Nature Index~\cite{nature-index-2026} and runs a decade-old ministry program for AI in science; the West has the Genesis Mission, Europe's factories and national machines; both sides underprovision by about an order of magnitude \full{24}. On the Chinese side that programme is older, more explicit and more industrial than the Western reader expects, and it is legible in primary policy text rather than inferred from behaviour; Chapter~\ref{c:china-ai4s} reads the record, and reads the parts of it that cut against this book.

\section{The bill: how much compute a nation's science needs}

\textbf{One labelling discipline governs every demand share below, and it has to be stated before any percentage is printed}, because the same physical capability---serving LLM tokens---carries entirely different weight in three regimes. \emph{Regime~G}, general and consumer AI reaching users through phones and applications, is essentially all token generation; the services its agents call are public and shared, not privately provisioned high-end compute. \emph{Regime~E}, enterprise streamlining---accounting, personnel, approvals, document flow---is likewise essentially all token generation: private data, substantial integration, the same serving capability. The United States and China contest G and E head to head, and every price table in Chapter~\ref{c:elements} is a statement about those two. \emph{Regime~S}, AI for science and industry, is this chapter's subject and is not built the same way: it is dominated by \emph{non-LLM} workloads---simulation, training of models that are not language models, data processing at instrument rate, integration with physical apparatus---with a token layer for agent orchestration that is a minority of the requirement and, on the scenario below, shrinking. Nothing in this book's numbers is contradictory once this is labelled, and much of the public argument about AI infrastructure is incoherent because it is not: ``language models are one seventh of the workload'' is a Regime~S statement and is evidence about nothing in G or E; ``token serving is the entire business'' is a G/E statement and is evidence about nothing in S. \textbf{Every share in this chapter is a share of Regime~S demand unless it says otherwise.}

Sovereignty becomes a checkable number only if the demand side is measured. National AI-for-science infrastructure has been sized by budget envelopes, vendor availability and flagship extrapolation; the full edition follows a measurement-anchored framework developed at RIKEN that replaces those with a vendor-neutral unit and a two-layer demand model~\cite{ai4s-framework} \full{26}. A requirement denominated in BGPU is procurable from any vendor, which is what a public planning framework must guarantee and a vendor-named unit cannot. The demand anchor is a measured 56-day campaign on a GB200-class partition: 101 users, 53,000 jobs, 340,000 GPU-hours, the partition saturated within five weeks~\cite{ai4s-framework}. The two layers are token service---researchers driving a trillion-parameter master model at million-token context that delegates to smaller sub-agents---and scientific execution. The structural finding is that ten percent of the tokens, generated at million-token context, cost nine tenths of the token-service GPUs, so national demand is overwhelmingly a long-context memory-capacity and bandwidth problem: the sizing framework and the architecture argument are the same argument in different units.

Three countries are instantiated, and the requirement tracks the researcher population closely because the per-researcher core sits near 0.94 BGPU: Japan needs roughly 497--510 thousand BGPU by 2030, the United States 1.09--1.17 million, Singapore 55--59 thousand, with a band of at least $\pm$40\% before local measurement, and the Singapore and US figures should not be quoted more precisely than transferred behavioural assumptions permit. Announced national infrastructures provision 0.02--0.10 BGPU per researcher against a modeled requirement near one---about ten times under for the United States and Japan, fifty times for Singapore. One side has the measurement without the money; the other has the money without the measurement. A nation that rents its master layer has rented the majority of its research computing; sovereignty has a computable price, about \$1.84--2.19 per delivered accelerator-hour, which makes pooling across time zones rational.

\emph{The gap, tested against October's money.} The \$2.4B of mostly in-kind Genesis Mission pledges announced on 8~October~\cite{sci-wh-summit} [V], valued at delivered cost within one year, buys at most 12--15\% of one year of the US Regime~S requirement [M]; Japan's request (Chapter~\ref{c:strategies}, where the author's interest in it is declared) targets tenfold shared compute by FY2030 but is unappropriated~\cite{gov-mext-fy2027} [P]. The measurement and the money met on neither side \full{26}.

\textbf{And the canonical serving configuration has an expiry date.} The framework prices its own obsolescence: on a DeepSeek-V4.1-Flash-class serving generation---aggressive sparsity, cache-resident long context, a per-token key--value footprint small enough that million-token master context stops costing what it costs today---cutting the serving coefficient by 3.25$\times$ takes Japan's token-service layer from 72.7 to 22.4 thousand BGPU while the national requirement falls only from 497 to 446 thousand, about 10\%, because the other six sevenths are untouched by a serving advance. The token layer's share of Regime~S demand falls from 14.3--14.6\% to about 5.0\% [F/S]. The published work driving that advance is Chinese. Both figures are carried forward rather than one replacing the other, because the 88.7\% master term is a property of a serving generation and not of agentic science. \emph{This strengthens the procurement argument rather than weakening it}: the commodity layer a middle power can simply buy shrinks from a seventh to a twentieth, while the layer nobody is organized to sell rises from six sevenths to nineteen twentieths. \emph{Against it}, a token 3.25 times cheaper is an invitation to run more agents at longer context, and the scenario holds behaviour fixed by construction; arithmetic cannot settle which way the layer moves.


\section{The third path: an international consortium}

For every country that is neither the United States nor China, the two available strategies are both unsatisfactory. Permanent dependence on a few foreign commercial providers cedes cost, control and continuity; isolated national programs duplicate the most expensive activity in the industry at sub-scale, at valuations that assume rents no one may collect. What is needed is a \emph{global} form of sovereign capability---not national isolation but shared ownership of the knowledge, data practices, compute access and institutional capacity required to build and operate scientific AI~\cite{consortium-record} \full{27}.

\emph{The consortium's emphasis was never competing with China---or with anyone---for widespread LLM use}, and so the shrinking token layer does not shrink the case for it. Nothing here proposes that a coalition of middle powers out-serve the two principals in Regimes~G and~E, where the contest is head to head and price is already falling toward marginal cost without assistance. The justifications are sovereignty, access, shared development, data stewardship, resilience and the avoidance of duplicated sub-scale effort, and none of the six has a token term that a cheaper serving generation removes; several get cheaper to exercise. The effort may also overlap and converge with general-purpose open frontier work such as NVIDIA's Nemotron line rather than competing with it. What follows is the financial argument, which is the one usually made and the one most often made backwards.

The decisive argument is financial. In a realistic AI-for-science portfolio, frontier language models and agents are roughly 14\% of the physical workload; simulation and domain-specific AI take the rest~\cite{consortium-record}. That figure is usually deployed as reassurance, and it is the opposite once the price asymmetry is applied: purchased frontier tokens carry an estimated 20--60-fold premium over equivalent service from institutionally controlled infrastructure running open models [A/M], so a Regime~S token-layer share of one seventh becomes 77.4\% of the budget at a twenty-fold premium and 91.1\% at sixty, and 33.1--54.1\% if only a tenth of the layer is escalated (Figure~\ref{fig:ai4s-cost}); on the 3.25$\times$ serving scenario above, with the layer at about a twentieth, the same two premia give 48.5\% and 73.8\%---smaller, and still the majority of the budget. Whoever sets the frontier token price sets the science budget of every institution that buys it: proposition P3 applied to national research capacity. This is also why the Chinese price collapse is strategically effective without coercion---at a premium near unity, which cheap open weights on domestic hardware deliver, the budget problem dissolves, and the institution that solved it that way has acquired a dependency it did not negotiate.

\begin{figure}[htbp]
\centering
\begin{tikzpicture}
\begin{axis}[bookaxis, height=5.4cm,
  xlabel={Price premium $\rho$ of purchased frontier tokens over self-operated open service},
  ylabel={Share of total expenditure (\%)},
  title={A minority workload becomes the majority of the budget},
  xmin=1, xmax=60, ymin=0, ymax=100,
  xtick={1,10,20,30,40,50,60},
  legend style={at={(0.98,0.32)}, anchor=east},
]
\addplot[color=sOne, mark=none, samples=80, domain=1:60] {100*0.14*x/(0.86+0.14*x)};
\addlegendentry{$f=14\%$ (measured workload share)}
\addplot[color=sThree, mark=none, samples=80, domain=1:60, densely dashed] {100*0.05*x/(0.95+0.05*x)};
\addlegendentry{$f=5\%$}
\addplot[color=sFour, mark=none, samples=80, domain=1:60, densely dotted] {100*0.25*x/(0.75+0.25*x)};
\addlegendentry{$f=25\%$}
\addplot[only marks, mark=*, mark size=2.4pt, color=sTwo] coordinates {(20,76.5) (60,90.7)};
\node[font=\scriptsize\color{inkSec}, anchor=south east] at (axis cs:20,78) {76.5\%};
\node[font=\scriptsize\color{inkSec}, anchor=south east] at (axis cs:59,92) {90.7\%};
\end{axis}
\end{tikzpicture}
\caption[The AI4S expenditure asymmetry]{Why the token price is the science budget. With a fraction $f$ of
physical work served by purchased frontier tokens at relative unit price $\rho$, the share of total expenditure that
work consumes is $\Sigma_{\mathrm{front}}=f_{\mathrm{front}}\rho/((1-f_{\mathrm{front}})+f_{\mathrm{front}}\rho)$. At the measured $f_{\mathrm{front}}\approx14\%$ and the estimated 20--60$\times$ premium, a
one-seventh workload share becomes 76.5--90.7\% of the budget [A/M]. The curve's shape,
not its exact position, is the policy argument: at $r$ near unity---which cheap open weights on owned hardware
deliver---the problem dissolves, and the institution that solved it that way has acquired a dependency it never
negotiated.}
\label{fig:ai4s-cost}
\end{figure}


Open pretraining is therefore insurance, not a benchmark contest: a consortium's own open model may never be the best production choice, and building it is still correct, because the return on pretraining is people who know how to do it. Near-term science should meanwhile use the best open-weight models, including Chinese ones, without embarrassment. The insurance covers three risks, none presuming bad faith---continuity, verification (Chapter~\ref{c:trust}), and regulation in dual-use science---and Chapter~\ref{c:china-ai4s} finds each of the three instantiated rather than hypothetical. Four functions define the institution: federated compute under common accounting, its first proof of concept the joint training of a fifteen-trillion-token open model split across compatible systems in several countries; scientific data stewardship, because science is the one community for which provenance and reproducibility are already requirements; model production with lifecycle obligations, since a clinical model may take three years to certify and must be supportable five years later---the strongest reason this cannot be a rotating academic project; and professional operation, because convening is cheap and operating is not. Industry participates without expropriation, on an open accountable base with closed private layers above it, and the Global South is treated as talent, questions and data rather than as aid, because the alternative is not neutrality but ceding the default to whoever shows up with a stack.

\emph{Pledges are not operated capability}: credits end with the pledge, model access is bought at the premium, and \$1.5B of the \$2.4B comes from the two accelerator vendors whose devices the mission's flagship systems use~\cite{sci-cdo-genesis} [A].

The organizational lesson from China is not a command structure but the reverse: government sets the stage and allows intense bottom-up activity, so a consortium funding one model team in one country reproduces the failure mode rather than the success. Table~\ref{tab:consortium} gives ten operational design parameters and the failure mode each avoids, and the full edition's feasibility appendix finds that of fourteen legal and engineering problems a founding instrument must solve, nine have answers on the shelf, three require legal invention, and none is a reason to wait \fullapp{F}. The governing framing is neutrality, because a certificate issued by one bloc about the other bloc's artifacts is worth nothing to anybody. Sequence: frame; prove federation with three to five committing countries and one visible cross-border pilot; then institutionalize. A five-country story with a working pilot beats a large inactive membership, and a proposal that names functions but not instruments is a wish.

\begin{table}[htbp]
\centering\footnotesize
\caption{The consortium, specified operationally: six of the ten design parameters the full edition states, with the failure mode each avoids [S].}
\label{tab:consortium}
\begin{tabular}{@{}p{32mm}p{60mm}p{56mm}@{}}
\toprule
\textbf{Parameter} & \textbf{Recommended answer} & \textbf{Failure mode if deferred} \\
\midrule
Legal form, host & Treaty-adjacent foundation in a neutral, non-aligned-capable jurisdiction & Seen as a bloc instrument; loses the Global South and one principal \\
Membership, votes & Separate classes for states, centres, universities, firms; contribution-weighted votes capped below a blocking share & Largest contributor becomes owner \\
Export-control compliance & Per-jurisdiction architecture, segmented enclaves & One member's rule governs everyone \\
Compute allocation & BGPU unit, peer review, reserved and burst tiers, cross-border settlement & Federation degrades to a mailing list \\
Certification authority & Operates the Level 0--5 ladder, publishes irrespective of origin & Vacancy stays vacant, or is filled by a supplier \\
Chinese and US technology & Capability- and evidence-based; no origin rule either way & Forfeits neutrality \\
\bottomrule
\end{tabular}
\end{table}

This is where the book's threads terminate---the institutional form of ``take the commons, verify it, contribute to it, and own the layer that makes it safe to use.'' Whether it gets built in the next two years or the next ten is, on the evidence of the precedents, the difference between shaping the outcome and documenting it.
